\RequirePackage{pdfmanagement}
\SetKeys[document/metadata]{pdfversion=1.7, pdfstandard=A-2u, lang=en-GB}
\input{preamble/pdfA-compliance/glyphtounicode}
\documentclass[11pt,a4paper]{report}

\newcommand*{\PaperTitle}{The discrete critical coagulation--fragmentation conjecture}
\newcommand*{\PaperTitlePlain}{The discrete critical coagulation-fragmentation conjecture}
\newcommand*{\PaperShortTitle}{Discrete critical coagulation--fragmentation}
\newcommand*{\PaperSubtitle}{chemistry-open-problems \textperiodcentered{} 20 October 2026 \textperiodcentered{} independently verified, not refereed}
\newcommand*{\PaperCredits}{Part (ii), the case m = 1, part (i) as a corollary, final manuscript: Claude Opus 5.5 (effort xhigh)\\
Novelty audit, stationary singularity analysis, first draft: GPT-6-Astra (effort xhigh)\\
Independent verification: Claude Opus 5.5 (effort xhigh)}
\newcommand*{\PaperKeywords}{coagulation, fragmentation, gelation, mass conservation, Hamilton-Jacobi equations}
\input{chemop-preamble}
\addbibresource{references.bib}
\newcommand{\NN}{\mathbb N}
\newcommand{\eps}{\varepsilon}
\newcommand{\rtil}{\tilde\rho}

\begin{document}
\makepapertitle

\begin{paperabstract}
Jang and Tran conjectured that the discrete coagulation--fragmentation equation with multiplicative
coagulation kernel \(ij\) and constant fragmentation kernel has, in the critical mass window
\(1/2\le m<1\), a unique stationary state of mass \(m\) given by an explicit recursion, and a unique global
mass-conserving classical solution. We prove both statements, for every mass \(0<m<1\), and the
time-dependent one also for \(m=1\); the initial second moment assumption is not needed, only finite mass.
With Jang and Tran's theorem that no global mass-conserving solution exists for \(m>1\), this settles the
discrete critical-mass conjecture, gelation for \(m>1\) being understood in that sense.
The zeroth-order coefficient of the associated Hamilton--Jacobi equation is \(1/\tau(x)\) with
\(\tau(x)=2\tanh(x/2)\). With \(x\) replaced by \(\tau\), Van's continuous half-slope invariant has a discrete
analogue that propagates along characteristics with a nonnegative source. It yields a curvature barrier for
all \(m\le1\), hence a global classical solution of the Hamilton--Jacobi equation by characteristics, with
slope \(m\) at the boundary. Comparison with positive finite truncations then excludes mass loss; no
viscosity-solution theory is used. Once mass is conserved the system is lower triangular, which gives
uniqueness, convergence to the stationary state, and its positivity and mass identity. The stationary
coefficients decay like \(c\,R^{-n}n^{-5/2}\) with \(R>1\). The proofs have been checked by an independent
verifier; they have not been refereed.
\end{paperabstract}

\section*{The question}

The source is Jang and Tran \autocite{JangTran2025}. The system is, for cluster densities
\(\rho_n=\rho(n,t)\ge0\), \(n\in\NN\),
\begin{equation}\label{eq:system}
\dot\rho_n=\frac12\sum_{i=1}^{n-1}i(n-i)\rho_i\rho_{n-i}-n\rho_n\sum_{j\ge1}j\rho_j
-\frac{n-1}{2}\rho_n+\sum_{j>n}\rho_j ,
\end{equation}
which is their (1.1) with kernels (1.2), \(a(i,j)=ij\), \(b(i,j)=1\). Their assumption (1.3) is
\(\sum_n n^2\rho_0(n)<\infty\), and their recursion (1.9) is
\begin{equation}\label{eq:recursion}
\rtil_n=\frac{2m(1-m)+\sum_{i=1}^{n-1}i(n-i)\rtil_i\rtil_{n-i}-2\sum_{i=1}^{n-1}\rtil_i}{(2m+1)n+1},
\qquad n\ge1 .
\end{equation}
Conjecture~2 of arXiv:2409.17974v1 (printed p.~7, end of Section~1.4), with notation typeset equivalently:
\begin{openquote}
Conjecture 2. Consider (1.1) with kernels (1.2) and initial data satisfying (1.3).
(i) For \(1/2<m<1\), (1.1) admits a unique stationary solution \(\tilde\rho:\NN\to[0,\infty)\) such that
\(m_1(\tilde\rho)=m\). Further, \(\tilde\rho\) satisfies the recursive formula (1.9).
(ii) For \(1/2\le m_1(\rho_0)<1\), (1.1) admits a unique mass-conserving solution \(\rho\) for all time.
This solution \(\rho\) is in fact a classical solution.
\end{openquote}
Jang and Tran proved existence for \(m<1/2\) under (1.3) (their Theorem~1.3), stationary states for
\(m\le1/2\) (Lemma~5.7), uniqueness of mass-conserving solutions for \(m\le1\) (Proposition~1.2) and
non-existence of global mass-conserving solutions for \(m>1\) (Theorem~1.1). Their Conjecture~1, which
they attribute to \autocite{EMP2002,ELMP2003}, asserts, for initial data satisfying (1.3), that \(m=1\) is
the critical mass: gelation for \(m>1\), a unique global mass-conserving solution for \(0<m\le1\).
For continuous sizes, Tran and Van \autocite{TranVan2022} proved existence for \(m<1/2\), and
Van's preprint \autocite{Van2026} extends existence to the window \(1/2\le m\le1\); Van writes, in
Section~1.2, p.~4: ``so the discrete analogue of the window \(1/2\le m\le1\) remains open as of this writing.''
The novelty audit is recorded in \texttt{references.md}; no earlier solution was found.

\paragraph{Notation and solution concepts.}
\(m_k(\rho)=\sum_n n^k\rho_n\); \(K(\rho)_n\) is the right side of \eqref{eq:system}, with the actual first
moment in the coagulation loss; \(\ell^1_1=\{u:\sum n|u_n|<\infty\}\). For bounded \(g:\NN\to\mathbb R\),
Jang and Tran's weak form (1.5) reads
\begin{equation}\label{eq:weak}
\frac{d}{dt}\sum_n g(n)\rho_n=\frac12\sum_{j,k}\bigl(g(j{+}k)-g(j)-g(k)\bigr)jk\rho_j\rho_k
-\frac12\sum_j\sum_{k=1}^{j-1}\bigl(g(j)-g(k)-g(j{-}k)\bigr)\rho_j .
\end{equation}
Their Definition~1 calls \(\rho\) a solution if \eqref{eq:weak} holds for every bounded \(g\), without fixing
the time regularity. We call \(\rho\) a \emph{mass-conserving solution} with datum \(\rho^0\) if
\(\rho(t)\ge0\), \(m_1(\rho(t))=m_1(\rho^0)\) for all \(t\ge0\), and for every bounded \(g\) the function
\(t\mapsto\sum g(n)\rho_n(t)\) is absolutely continuous on compact intervals and satisfies \eqref{eq:weak}
for almost every \(t\). This is the weakest reading; uniqueness is proved in this class. It is
\emph{classical} if every \(\rho_n\) is \(C^1\) on \([0,\infty)\) and \eqref{eq:system} holds for every
\(t\ge0\). A \emph{stationary solution} is a time-independent mass-conserving solution, that is,
\(\tilde\rho\ge0\), \(m_1(\tilde\rho)<\infty\), \(K(\tilde\rho)=0\).

\begin{theorem}[global well-posedness]\label{thm:evolution}
Let \(\rho^0\ge0\) with \(m=\sum_n n\rho^0_n\in(0,1]\); no other moment is assumed. Then \eqref{eq:system}
has exactly one mass-conserving solution with datum \(\rho^0\). It is classical, every component is
\(C^\infty\) on \([0,\infty)\),
\(\rho\in C([0,\infty);\ell^1_1)\cap C^1([0,\infty);\ell^1)\), the weak form \eqref{eq:weak} holds with a
classical derivative at every \(t\ge0\) for every bounded \(g\), and
\[
\sum_n\rho_n(t)=m(1-m)+e^{-t/2}\Bigl(\sum_n\rho^0_n-m(1-m)\Bigr).
\]
\end{theorem}

\begin{theorem}[stationary states and relaxation]\label{thm:stationary}
Let \(0<m<1\). The numbers \eqref{eq:recursion} are strictly positive, with
\(\sum_n\rtil_n=m(1-m)\) and \(\sum_n n\rtil_n=m\); \(\rtil\) is a stationary solution and the only stationary
solution with \(m_1=m\). Every mass-conserving solution of mass \(m\) converges to \(\rtil\) in \(\ell^1_1\).
For \(m=1\), every component of the solution of \zcref{thm:evolution} tends to \(0\) while its mass stays
equal to \(1\).
\end{theorem}

\begin{theorem}[decay of the stationary state]\label{thm:asymptotics}
For \(0<m<1\) the generating function \(\sum_n\rtil_nz^n\) has radius \(R=R(m)\in(1,\infty)\), and
\[
\rtil_n=\frac{b}{2\sqrt\pi}\,R^{-n}n^{-5/2}\bigl(1+O(n^{-1})\bigr),\qquad
b^2=\frac{(R-1)^2(m+\frac12)+\frac32R}{R^2-1}.
\]
\end{theorem}

Thus Conjecture~2 holds: part (i) is \zcref{thm:stationary} and part (ii) is \zcref{thm:evolution}, both
for every \(0<m<1\), and part (ii) without (1.3). \zcref{thm:evolution} at \(m=1\), together with Jang and
Tran's Theorem~1.1 for \(m>1\), settles their Conjecture~1 in the discrete setting. Here ``gelation for
\(m>1\)'' is understood as in their Theorem~1.1: no solution conserves mass for all time. This also follows
from \zcref{prop:closed}, whose derivation of \(n_0(t)\) holds on any interval \([0,T]\) and makes \(n_0\)
vanish at \(T=2\log\bigl(1+n_0(0)/(m(m-1))\bigr)\), so no solution conserves mass on such an interval.
Whether global weak solutions that lose mass exist for \(m>1\) is not addressed here.
The status is \emph{solved}: the proofs were checked by an independent verifier (file
\texttt{verification.md}); they have not been refereed.

\paragraph{Strategy.}
Write \(F(x,t)=\sum_n(1-e^{-nx})\rho_n(t)\) (Jang and Tran's discrete Bernstein transform). The zeroth-order
coefficient of its Hamilton--Jacobi equation (Jang--Tran (1.7)) is
\(\frac12+\frac1{e^x-1}=\frac12\coth\frac x2=1/\tau(x)\), \(\tau(x)=2\tanh(x/2)\): it is Tran and Van's
continuous equation with \(x\) replaced by \(\tau(x)\) in the term \(F/x\).
(1) The quantity \(W=\tau(m+F_x)-2F\) is nonnegative for the transform of any nonnegative sequence of mass
\(m\) and propagates along characteristics with a nonnegative source; it is the discrete analogue of Van's
half-slope invariant \autocite{Van2026}. (2) It gives the curvature barrier \(\tau F_{xx}>-1\) for all
\(m\le1\); Jang and Tran's version of this barrier (their Lemma~3.6) needed \(m<1/2\). (3) Characteristics
leave the half-line through \(x=0\); with the barrier they never cross, so the Hamilton--Jacobi equation has a
global classical solution, with slope \(m\) at \(x=0\) by an ODE argument along characteristics. (4) Finite
truncations of \eqref{eq:system} are positive and conserve mass exactly; their transforms are
supersolutions; transport comparison puts them above the characteristic solution, so a limit point cannot
lose mass. (5) With constant mass, the number of clusters is explicit and \eqref{eq:system} becomes lower
triangular: uniqueness, relaxation and part (i) follow.

\section{The closed triangular system and uniqueness}\label{sec:closed}

\begin{proposition}\label{prop:closed}
Let \(\rho\) be a mass-conserving solution of mass \(m>0\) and \(n_0(t)=\sum_n\rho_n(t)\). Then
\(n_0(t)=m(1-m)+e^{-t/2}(n_0(0)-m(1-m))\); every \(\rho_n\) is \(C^\infty\) on \([0,\infty)\) and, with
\(c_n=mn+\frac{n+1}{2}\),
\begin{equation}\label{eq:closed}
\dot\rho_n+c_n\rho_n=n_0(t)-\sum_{j<n}\rho_j+\frac12\sum_{i=1}^{n-1}i(n-i)\rho_i\rho_{n-i}
\qquad(t\ge0);
\end{equation}
and two mass-conserving solutions with the same datum coincide.
\end{proposition}
\begin{proof}
All series converge absolutely because \(\sum n\rho_n=m\). With \(g=\mathbf 1_{\{n\}}\) in \eqref{eq:weak},
\(\rho_n\) is absolutely continuous and \(\dot\rho_n=K(\rho)_n\) almost everywhere, with coagulation loss
\(nm\rho_n\). With \(g\equiv1\) the right side of \eqref{eq:weak} is
\(-\frac{m^2}2+\frac12\sum(n-1)\rho_n=\frac12(m-m^2-n_0)\); this linear equation for the absolutely
continuous \(n_0\) gives the formula (Jang--Tran (1.6)). Since \(\sum_{j>n}\rho_j=n_0-\sum_{j\le n}\rho_j\),
\eqref{eq:closed} holds almost everywhere. Its right side involves only the smooth \(n_0\) and
\(\rho_1,\dots,\rho_{n-1}\). By induction on \(n\), \(\rho_n\) equals its variation-of-constants formula with
a continuous, then \(C^\infty\), forcing; so \(\rho_n\in C^\infty\) and \eqref{eq:closed} holds everywhere.
Two solutions with the same datum have the same \(n_0\), and \eqref{eq:closed} for \(n=1,2,\dots\) are linear
scalar equations with identical forcings, by induction.
\end{proof}

Jang and Tran proved this uniqueness (their Proposition~1.2) through uniqueness of viscosity solutions; the
triangular structure gives it directly.

\section{The auxiliary Hamilton--Jacobi equation}\label{sec:HJ}

Throughout this section \(m\in(0,1]\) and \(\rho^0\ge0\) with \(\sum n\rho^0_n=m\) are fixed. Put
\[
\tau(x)=2\tanh\frac x2,\qquad k=\tau'=1-\frac{\tau^2}4\in(0,1],
\]
\[
H_m(x,u,p)=\frac12(p-m)(p-m-1)+\frac{u}{\tau(x)}-m,\qquad
F_0(x)=\sum_n(1-e^{-nx})\rho^0_n .
\]
For \(x>0\): \(0<\tau(x)\le x\), \(\tau<2\), and \(1/\tau(x)=\frac12+\frac1{e^x-1}\). The function \(F_0\) is
\(C^\infty\) on \((0,\infty)\), positive, with \(0\le F_0'=\sum ne^{-nx}\rho^0_n\le m\), \(F_0''\le0\) and
\(F_0(x)/x\uparrow m\) as \(x\downarrow0\) (monotone convergence). We study
\begin{equation}\label{eq:HJ}
F_t+H_m(x,F,F_x)=0\quad(x>0,\ t\ge0),\qquad F(\cdot,0)=F_0 .
\end{equation}
We do not assume that its solution is the transform of a nonnegative sequence: it only serves as a lower
barrier in Section~\ref{sec:existence}.

\begin{lemma}\label{lem:initial}
\(W_0:=\tau(m+F_0')-2F_0\ge0\) and \(-m/e\le\tau F_0''\le0\) on \((0,\infty)\).
\end{lemma}
\begin{proof}
With \(z=e^{-x}\), \(\tau=2(1-z)/(1+z)\). For an integer \(j\ge1\),
\((1+z)(1+z+\dots+z^{j-1})=1+2(z+\dots+z^{j-1})+z^j\), hence
\[
\tau j(1+z^j)-2(1-z^j)=\frac{2(1-z)}{1+z}\Bigl(j(1+z^j)-(1+z)\frac{1-z^j}{1-z}\Bigr)
\]
\[
=\frac{2(1-z)}{1+z}\sum_{i=1}^{j-1}(1-z^i)(1-z^{j-i})\ge0 .
\]
Multiply by \(\rho^0_j\) and sum, using \(\sum j\rho^0_j=m\). Next,
\(-\tau F_0''=\sum_j j\rho^0_j\,(\tau je^{-jx})\) and \(\tau je^{-jx}\le xje^{-jx}\le1/e\).
\end{proof}

\paragraph{Characteristics.}
For a label \(y>0\) let \((X,P,U)(t)=(X,P,U)(t;y)\) solve
\begin{equation}\label{eq:char}
\dot X=P-m-\tfrac12,\qquad
\dot P=\frac{k(X)}{\tau(X)^2}U-\frac{P}{\tau(X)},\qquad
\dot U=\tfrac12(P^2+m-m^2)-\frac{U}{\tau(X)},
\end{equation}
with \((X,P,U)(0)=(y,F_0'(y),F_0(y))\). These are
\(\dot X=\partial_pH_m\), \(\dot P=-\partial_xH_m-p\,\partial_uH_m\), \(\dot U=p\,\partial_pH_m-H_m\).
The vector field is \(C^\infty\) on \(\{X>0\}\times\mathbb R^2\); \([0,t_e(y))\) is the maximal interval of
existence. Put \(q=m-P\), \(v=q+\frac12\), \(W=\tau(X)(m+P)-2U\), \(D=U/\tau(X)-P\), \(\lambda=U/\tau(X)\),
\(\eta=m-\lambda\), \(\eps=\tau(X)^2v/4\). Direct differentiation gives the exact identities
\begin{gather}
\dot W+\frac{2+k}{2\tau}W=q(2m-q)(1-k),\qquad 1-k=\frac{\tau^2}4,\label{eq:W}\\
W=\tau(q-2D),\label{eq:WD}\\
\dot\lambda=\frac{(P-\lambda)^2+(m-\lambda)(1-m+\lambda)}{2\tau}-\frac{\eps\lambda}{\tau},\qquad
\dot\eta=-\frac{(\eta-q)^2+\eta(1-\eta)}{2\tau}+\frac{\eps\lambda}{\tau}.\label{eq:lambda}
\end{gather}
In the continuous case (\(\tau=x\), \(k=1\)) the source in \eqref{eq:W} vanishes and \eqref{eq:W} is Van's
equation (3.3); here it is nonnegative.

\begin{lemma}[invariant region]\label{lem:region}
For every \(y>0\) and \(t\in[0,t_e(y))\): \(0\le P\le m\), \(U>0\), \(W\ge0\). Consequently
\(0<U\le m\tau(X)<2m\), \(0\le\lambda\le\frac{m+P}2\le m\), \(-(m+\frac12)\le\dot X\le-\frac12\),
\(y/(m+\frac12)\le t_e(y)\le2y\), and \(X(t)\downarrow0\) as \(t\uparrow t_e(y)\).
\end{lemma}
\begin{proof}
Fix \(T<t_e(y)\); on \([0,T]\), \(X\ge\delta>0\) and \(|P|,|U|\le A\). Let \(\delta'=\tau(\delta)\) and
\(V=P^-+(P-m)^++U^-+W^-\), where \(a^\pm=\max(\pm a,0)\). \(V\) is Lipschitz and almost everywhere
\(V'=-\dot P\mathbf 1_{P<0}+\dot P\mathbf 1_{P>m}-\dot U\mathbf 1_{U<0}-\dot W\mathbf 1_{W<0}\). On
\(\{P<0\}\), \(-\dot P=-kU/\tau^2+P/\tau\le U^-/\delta'^2\). On \(\{P>m\}\), since
\(2U=\tau(m+P)-W\le\tau(m+P)+W^-\) and \(k\le1\),
\(\dot P\le\frac{k(m+P)-2P}{2\tau}+\frac{W^-}{2\tau^2}\le\frac{m-P}{2\tau}+\frac{W^-}{2\delta'^2}\le\frac{W^-}{2\delta'^2}\).
On \(\{U<0\}\), \(-\dot U=-\frac12(P^2+m(1-m))+U/\tau<0\); this uses \(m\le1\). On \(\{W<0\}\),
\eqref{eq:W} gives \(-\dot W\le-q(2m-q)(1-k)\), and \(q(2m-q)\ge0\) if \(0\le P\le m\),
\(q(2m-q)=m^2-(P^-)^2\ge-AP^-\) if \(P<0\), \(q(2m-q)=-(P-m)^+(2m+(P-m)^+)\ge-(2m+A)(P-m)^+\) if \(P>m\).
Hence \(V'\le LV\) almost everywhere for a constant \(L\), and \(V(0)=0\) by \zcref{lem:initial}. Grönwall gives
\(V\equiv0\). Then \(U\le\tau(m+P)/2\le m\tau\) and \(\lambda\le(m+P)/2\). From \(\dot U\ge-U/\tau\),
\(U(t)\ge F_0(y)\exp(-\int_0^t ds/\tau(X))>0\). The bounds on \(\dot X\) give the bounds on \(t_e\). If \(X\)
stayed above some \(\delta>0\), the solution would remain in a compact subset of the domain and could be
continued; as \(X\) decreases, \(X\to0\).
\end{proof}

\begin{lemma}[strip condition and curvature barrier]\label{lem:barrier}
Let \(J=\partial_yX\), \(\Pi=\partial_yP\), \(\Upsilon=\partial_yU\) on
\(\Omega=\{(y,t):y>0,\ 0\le t<t_e(y)\}\). Then \(\Upsilon=PJ\) and \(J>0\) on \(\Omega\), and
\(\Gamma:=\tau(X)\Pi/J\) satisfies \(-1<\Gamma\le\sqrt{2m}\).
\end{lemma}
\begin{proof}
Differentiating \eqref{eq:char} in \(y\), \(A=\Upsilon-PJ\) satisfies \(\dot A=-\partial_uH_m\,A=-A/\tau\),
and \(A(0)=F_0'(y)-F_0'(y)=0\); so \(A\equiv0\). With \(\kappa=1/\tau\) this gives \(\dot J=\Pi\) and
\(\dot\Pi=-J(\kappa''U+2\kappa'P)-\kappa\Pi\). Let \(t_1=\sup\{t<t_e:J>0\text{ on }[0,t]\}\). On \([0,t_1)\),
\(r=\Pi/J\) satisfies \(\dot r=-r^2-\kappa r-2\kappa'P-\kappa''U\), and \(\Gamma=\tau r\) satisfies
\begin{equation}\label{eq:Gamma}
\tau\dot\Gamma=-\Gamma^2-\Gamma\bigl(1+k(q+\tfrac12)\bigr)-2kD .
\end{equation}
By \eqref{eq:WD} and \zcref{lem:region}, \(D\le q/2\); also \(D\ge-P\ge-m\). Initially
\(\Gamma(0)=\tau(y)F_0''(y)\in[-m/e,0]\). At a first time with \(\Gamma=-1\), \eqref{eq:Gamma} gives
\(\tau\dot\Gamma=k(q+\frac12-2D)\ge k/2>0\), which is impossible; at a first time with
\(\Gamma=\sqrt{2m}\) it gives \(\tau\dot\Gamma\le-2m-\sqrt{2m}+2km<0\), also impossible. So
\(-1<\Gamma\le\sqrt{2m}\) on \([0,t_1)\). If \(t_1<t_e\), then \(X\ge X(t_1)>0\) on \([0,t_1]\), so
\(r\ge-1/\tau(X(t_1))\) and \(J(s)=\exp\int_0^sr\ge\exp(-t_1/\tau(X(t_1)))\) on \([0,t_1)\); by continuity
\(J>0\) on a neighbourhood of \(t_1\), a contradiction. Hence \(J>0\) on \(\Omega\).
\end{proof}

With only the bound \(D\le q\) (which is what \(F\le m(1-e^{-x})\) gives) the right side of \eqref{eq:Gamma}
at \(\Gamma=-1\) is only known to be at least \(k(\frac12-q)\), positive when \(q<\frac12\): this is the
origin of the restriction \(m<1/2\) in Jang and Tran's Lemma~3.6. The half-slope invariant is used exactly here.

\begin{proposition}[global classical solution]\label{prop:global}
For each \(t\ge0\), \(\Omega_t=\{y:t_e(y)>t\}\) is an interval \((a(t),\infty)\), with \(a(0)=0\) and
\(a(t)\ge t/2\), and \(X(t;\cdot)\) maps it increasingly onto \((0,\infty)\). The map
\(\Phi(y,t)=(X(t;y),t)\) is a \(C^\infty\) diffeomorphism of \(\Omega\) onto \((0,\infty)\times[0,\infty)\).
The function \(F=U\circ\Phi^{-1}\) is \(C^\infty\) on \((0,\infty)\times[0,\infty)\) (one-sided derivatives
at \(t=0\)), \(F(\cdot,0)=F_0\), \(F_x=P\circ\Phi^{-1}\), \(F\) solves \eqref{eq:HJ}, and
\[
0\le F_x\le m,\qquad 0<F\le m\tau(x),\qquad \tau(m+F_x)-2F\ge0,\qquad -1<\tau F_{xx}\le\sqrt{2m}.
\]
\end{proposition}
\begin{proof}
\(\Omega\) is relatively open (lower semicontinuity of \(t_e\)) and \(\Phi\) is \(C^\infty\) with Jacobian
determinant \(J>0\). \emph{Interval property.} Let \(y<y'\), \(y\in\Omega_t\), and
\(s^*=\sup\{s\le t:[y,y']\times[0,s]\subset\Omega\}\). For \(s<s^*\) and \(\eta\in[y,y']\),
\(X(s;\eta)\ge X(s;y)\ge X(t;y)>0\), because \(J>0\) and \(X\) decreases in time. So on \([0,s^*)\) these
characteristics stay in the compact set \([X(t;y),y']\times[0,m]\times[0,2m]\), near which the vector field is
smooth; a uniform extension time shows \(s^*=t\) and \(y'\in\Omega_t\). Thus \(\Omega_t=(a(t),\infty)\). As
\(t_e(y)\le2y\), \(a(t)\ge t/2\). \emph{Onto.} \(X(t;y)\ge y-(m+\frac12)t\to\infty\). If \(X(t;y)\ge\ell>0\) on
\(\Omega_t\), then \(a=a(t)>0\), \(t_e(a)\le t\), and continuous dependence gives, for \(s<t_e(a)\),
\(X(s;a)=\lim_{y\downarrow a}X(s;y)\ge\ell\), contradicting \(X(s;a)\to0\) (\zcref{lem:region}). So
\(\Phi\) is a bijection with positive Jacobian, hence a diffeomorphism. \emph{Equation.} Differentiating
\(U(t;y)=F(X(t;y),t)\) in \(y\) gives \(\Upsilon=F_xJ\), so \(F_x=P\) by \zcref{lem:barrier}; in \(t\),
\(\dot U=F_x\dot X+F_t\), so \(F_t=\dot U-P\partial_pH_m=-H_m\). Differentiating \(F_x(X,t)=P\) in \(y\)
gives \(F_{xx}=\Pi/J\). The bounds are \zcref{lem:region,lem:barrier}.
\end{proof}

\begin{lemma}[comparison]\label{lem:comparison}
Let \(G\in C^1((0,\infty)\times[0,\infty))\) satisfy \(G_t+H_m(x,G,G_x)\ge0\), \(0\le G_x\le m\) and
\(G(\cdot,0)\ge F_0\). Then \(G\ge F\).
\end{lemma}
\begin{proof}
Since \(H_m(x,u,p)-H_m(x,\tilde u,\tilde p)=(p-\tilde p)\bigl(\frac{p+\tilde p}2-m-\frac12\bigr)+(u-\tilde u)/\tau\),
\(w=F-G\) satisfies \(w_t+bw_x+w/\tau\le0\) with \(b=\frac{F_x+G_x}2-m-\frac12\in[-(m+\frac12),-\frac12]\)
continuous. Fix \(x_1,t_1>0\) and let \(Y\) solve \(\dot Y=b(Y,t)\), \(Y(t_1)=x_1\), backwards in time
(Peano). Backwards, \(Y\) increases: \(Y(t)\in[x_1+\frac{t_1-t}2,x_1+(m+\frac12)(t_1-t)]\), so \(Y\) exists
on \([0,t_1]\) inside \((0,\infty)\). Then \(g(t)=w(Y(t),t)\) satisfies \(\dot g\le-g/\tau(Y)\), so
\(w(x_1,t_1)\le w(Y(0),0)\exp(-\int_0^{t_1}1/\tau(Y))\le0\). No boundary condition at \(x=0\) is needed.
\end{proof}

\begin{lemma}[slope at the boundary]\label{lem:slope}
(a) For \(m\in(0,1]\) and every \(t\ge0\), \(F(x,t)/x\to m\) as \(x\downarrow0\).
(b) If \(m<1\), \(t_0>0\) and \(0<\delta<\min\{1,t_0/2\}\), then for \(0<x<\delta\) and \(t\ge t_0\)
\[
0\le m-\frac{F(x,t)}{\tau(x)}\le Kx^\alpha,
\]
\[
\alpha=\frac{1-m}{2m+1},\qquad K=m\delta^{-\alpha}+\frac{2C}{2-\alpha}\delta^{2-\alpha},\qquad
C=\frac{m(m+\frac12)}4 .
\]
\end{lemma}
\begin{proof}
At \(t=0\), (a) is \(F_0(x)/x\to m\). Let \(t>0\), \(0<\delta<\min\{1,t/2\}\), \(0<x<\delta\), and follow the
characteristic through \((x,t)\). Its label is \(y\ge x+t/2>\delta\) and \(X\) decreases at speed at least
\(\frac12\), so \(X(s_1)=\delta\) for some \(s_1\in[t-2\delta,t)\subset(0,t)\), and \(X\in[x,\delta]\) on
\([s_1,t]\). Regard \(\lambda,\eta\) on this segment as functions of \(X\); \(d/dX=-v^{-1}d/ds\),
\(\frac12\le v\le m+\frac12\le\frac32\), \(\eps\le\frac38X^2\), and \(\eps\lambda/\tau\le CX\).

(b) As \(\eta\le m<1\), \(\eta(1-\eta)\ge(1-m)\eta\), and \eqref{eq:lambda} gives
\(\dot\eta\le-\frac{(1-m)\eta}{2\tau}+CX\). Since \(\eta\ge0\) and \(\tau\le X\),
\(\frac{d\eta}{dX}\ge\frac{(1-m)\eta}{2v\tau}-\frac{CX}{v}\ge\frac{\alpha}{X}\eta-2CX\). Integrating
\((X^{-\alpha}\eta)'\ge-2CX^{1-\alpha}\) over \([x,\delta]\) with \(\eta(\delta)\le m\) gives the bound, whose
constant does not depend on \(t\ge t_0\).

(a) Let \(m\le1\) and \(\theta\in(0,m)\); choose \(\delta\) also with \(\frac38\delta^2\le\frac\theta4\). As
\(1-m+\lambda\ge\lambda\), \eqref{eq:lambda} gives \(\dot\lambda\ge\frac\lambda\tau(\frac{m-\lambda}2-\eps)\).
So on \([s_1,t]\), \(\dot\lambda\ge\theta\lambda/(4\tau)>0\) whenever \(\lambda\le m-\theta\): the set
\(\{\lambda\ge m-\theta\}\) is never left once entered, and while \(\lambda\le m-\theta\),
\(d\lambda/dX\le-\frac\theta6\lambda/X\), i.e. \(\lambda(X)\ge\lambda(\delta)(\delta/X)^{\theta/6}\). The entry
value satisfies \(\lambda(\delta)=F(\delta,s_1)/\tau(\delta)\ge\lambda_*:=\min_{s\in[0,t]}F(\delta,s)/\tau(\delta)>0\).
Hence \(F(x,t)/\tau(x)\ge m-\theta\) for \(x<\delta(\lambda_*/(m-\theta))^{6/\theta}\), while
\(F/\tau\le m\). Since \(\theta\) is arbitrary and \(\tau(x)/x\to1\), (a) follows.
\end{proof}

\section{Existence: positive truncations and comparison}\label{sec:existence}

For \(L\ge2\) the truncated system keeps the coagulation events with \(i+j\le L\):
\begin{equation}\label{eq:trunc}
\dot\rho^L_n=\frac12\sum_{i=1}^{n-1}i(n-i)\rho^L_i\rho^L_{n-i}-n\rho^L_n\sum_{j=1}^{L-n}j\rho^L_j
-\frac{n-1}{2}\rho^L_n+\sum_{j=n+1}^{L}\rho^L_j,\qquad1\le n\le L,
\end{equation}
with \(\rho^L_n(0)=\rho^0_n\) for \(n<L\) and \(\rho^L_L(0)=\rho^0_L+L^{-1}\sum_{j>L}j\rho^0_j\). Let
\(F^L(x,t)=\sum_{n\le L}(1-e^{-nx})\rho^L_n(t)\).

\begin{lemma}\label{lem:trunc}
\eqref{eq:trunc} has a unique global solution; it is nonnegative and \(\sum_{n\le L}n\rho^L_n\equiv m\).
Moreover \(F^L\) is \(C^1\), \(0\le F^L_x\le m\), \(F^L(\cdot,0)\ge F_0\), and
\begin{equation}\label{eq:super}
F^L_t+H_m(x,F^L,F^L_x)=E_L:=\frac12\sum_{\substack{i,j\le L\\ i+j>L}}ij\rho^L_i\rho^L_j(1-e^{-ix})(1-e^{-jx})\ \ge0 .
\end{equation}
\end{lemma}
\begin{proof}
The right side of \eqref{eq:trunc} is polynomial and quasi-positive (\(\rho^L\ge0\), \(\rho^L_n=0\) imply
\(\dot\rho^L_n\ge0\)), so the nonnegative orthant is invariant. Coagulation conserves \(i+j\) and fragmentation
conserves size, so \(\sum n\rho^L_n\) is constant, equal to \(\sum_{n<L}n\rho^0_n+L\rho^0_L+\sum_{j>L}j\rho^0_j=m\);
the solution is bounded, hence global. For \eqref{eq:super} take \(g(n)=1-z^n\), \(z=e^{-x}\), in the weak
form of \eqref{eq:trunc}. As \(g(i+j)-g(i)-g(j)=-(1-z^i)(1-z^j)\), coagulation contributes
\(-\frac12(\sum_nn(1-z^n)\rho^L_n)^2+E_L=-\frac12(m-F^L_x)^2+E_L\). For fragmentation,
\(\sum_{k=1}^{n-1}(g(n)-g(k)-g(n-k))=-(n-1)(1+z^n)+2(z-z^n)/(1-z)\); summing against \(-\frac12\rho^L_n\) with
\(\sum z^n\rho^L_n=n_0-F^L\), \(\sum nz^n\rho^L_n=F^L_x\) gives
\(\frac12(m+F^L_x)-F^L\bigl(\frac1{1-z}-\frac12\bigr)=\frac12(m+F^L_x)-F^L/\tau\). Hence
\(F^L_t=-\frac12(m-F^L_x)^2+\frac12(m+F^L_x)-F^L/\tau+E_L\), which is \eqref{eq:super}. The same computation
with \(E=0\) gives Jang--Tran's (1.7) for a mass-conserving solution. Finally
\(F^L(x,0)-F_0(x)=\sum_{j>L}j\rho^0_j\bigl(\frac{1-e^{-Lx}}L-\frac{1-e^{-jx}}j\bigr)\ge0\), since
\(s\mapsto(1-e^{-sx})/s\) decreases.
\end{proof}

\begin{proof}[of \zcref{thm:evolution}]
Uniqueness and the formula for \(n_0\) are \zcref{prop:closed}.

\emph{Barrier.} By \zcref{lem:trunc,lem:comparison}, \(F^L\ge F\) for every \(L\).

\emph{Compactness.} Extend \(\rho^L_n=0\) for \(n>L\). Then \(0\le\rho^L_n\le m/n\) and
\(|\dot\rho^L_n|\le3m^2/2+3m/2\). By Arzelà--Ascoli and a diagonal argument, along some \(L_k\to\infty\),
\(\rho^{L_k}_n\to\rho_n\) locally uniformly on \([0,\infty)\) for every \(n\), with \(\rho_n\ge0\) continuous,
\(\rho_n(0)=\rho^0_n\), and \(M(t):=\sum n\rho_n(t)\le m\) by Fatou. As \(\sum_{n>N}\rho^L_n\le m/(N+1)\),
\(F^{L_k}(x,t)\to F^\rho(x,t):=\sum(1-e^{-nx})\rho_n(t)\).

\emph{No mass loss.} \(F^\rho\ge F\), and \(1-e^{-nx}\le nx\) gives
\(M(t)\ge F^\rho(x,t)/x\ge F(x,t)/x\) for all \(x>0\). By \zcref{lem:slope}(a), \(M(t)\ge m\); so \(M(t)=m\)
for every \(t\ge0\).

\emph{Limit equation.} Fix \(t\). Coordinatewise convergence and equality of the masses give tightness: if
\(\sum_{n\le N}n\rho_n(t)>m-\epsilon\), then \(\sum_{n>N}n\rho^{L_k}_n(t)<2\epsilon\) for large \(k\). So
\(\mu^{L_k}_n(t)=\sum_{j\le L_k-n}j\rho^{L_k}_j(t)\to m\) for fixed \(n\). In the integral form of
\eqref{eq:trunc}, the coagulation gain is a finite sum, the loss \(n\rho^L_n\mu^L_n\le m^2\) converges
pointwise to \(nm\rho_n\) (dominated convergence), and the fragmentation gain converges uniformly by the
number-tail bound. Hence \(\rho_n(t)=\rho^0_n+\int_0^tK(\rho(s))_n\,ds\) for all \(n\) and \(t\).

\emph{Regularity.} If \(t_j\to t\), the nonnegative sequences \(\rho(t_j)\) converge coordinatewise to
\(\rho(t)\) with equal masses, so \(\sum n|\rho_n(t_j)-\rho_n(t)|=2\sum n(\rho_n(t)-\rho_n(t_j))^+\to0\)
(discrete Scheffé, dominated convergence): \(\rho\in C([0,\infty);\ell^1_1)\). The map
\(K:\ell^1_1\to\ell^1\) is locally Lipschitz (bilinear coagulation part of norm at most
\(\frac32\|u\|_{\ell^1_1}\|v\|_{\ell^1_1}\), linear fragmentation part of norm at most \(\frac32\)). So the integral
identity holds in \(\ell^1\) with continuous integrand, \(\rho\in C^1([0,\infty);\ell^1)\), and for bounded
\(g\), \(\sum g(n)\rho_n\) is \(C^1\) with derivative \(\sum g(n)K(\rho)_n\), the right side of \eqref{eq:weak}
after rearranging absolutely convergent double series. Thus \(\rho\) is a classical mass-conserving solution;
\zcref{prop:closed} gives \(C^\infty\) components. By uniqueness the whole family \(\rho^L\) converges.
\end{proof}

\begin{remark}
The transform of the solution is the characteristic solution: \(F^\rho=F\). Indeed \(F^\rho\) is \(C^1\) and
solves \eqref{eq:HJ} by the computation of \zcref{lem:trunc} with \(E=0\); \(w=F^\rho-F\) then solves
\(w_t+bw_x+w/\tau=0\), \(w(\cdot,0)=0\), and the argument of \zcref{lem:comparison} applied to \(\pm w\) gives
\(w=0\). This is not needed above.
\end{remark}

\section{Stationary states and relaxation}\label{sec:stationary}

\begin{proof}[of \zcref{thm:stationary}]
Let \(0<m\le1\) and \(\rho\) be the solution of \zcref{thm:evolution} for a datum of mass \(m\). In
\eqref{eq:closed}, \(n_0(t)\to m(1-m)\). If \(\dot u+cu=f(t)\) with \(c>0\) and \(f(t)\to f_\infty\), then
\(u(t)\to f_\infty/c\). Applied to \eqref{eq:closed} for \(n=1,2,\dots\) in turn, this gives
\(\rho_n(t)\to\rtil_n\), where \(c_n\rtil_n=m(1-m)-\sum_{j<n}\rtil_j+\frac12\sum i(n-i)\rtil_i\rtil_{n-i}\), which is
\eqref{eq:recursion} since \(2c_n=(2m+1)n+1\). Hence \(\rtil_n\ge0\). For \(m=1\) the recursion gives
\(\rtil\equiv0\).

Let now \(m<1\). By the number-tail bound, \(F^\rho(x,t)\to\tilde F(x):=\sum(1-e^{-nx})\rtil_n\) and
\(\sum\rtil_n=\lim n_0(t)=m(1-m)\). By the barrier step of the previous proof, \(F^\rho\ge F\), and
\zcref{lem:slope}(b) with \(t_0=1\) gives \(F^\rho(x,t)\ge\tau(x)(m-Kx^\alpha)\) for \(0<x<\delta\),
\(t\ge1\). Letting \(t\to\infty\), \(\tilde F(x)\ge\tau(x)(m-Kx^\alpha)\), so
\(\sum n\rtil_n=\lim_{x\downarrow0}\tilde F(x)/x\ge m\); Fatou gives the reverse inequality.
With these two sums, \(\sum_{j>n}\rtil_j=m(1-m)-\sum_{j\le n}\rtil_j\), and \eqref{eq:recursion} is the equation
\(K(\rtil)_n=0\): \(\rtil\) is stationary. Strict positivity: \(\rtil_1=m(1-m)/(m+1)>0\), and the numerator of
\eqref{eq:recursion} equals \(2\sum_{j\ge n}\rtil_j+\sum_{i=1}^{n-1}i(n-i)\rtil_i\rtil_{n-i}\), which is at least
\((n-1)\rtil_1\rtil_{n-1}>0\) if \(\rtil_{n-1}>0\). Uniqueness: a stationary \(\sigma\) with \(m_1(\sigma)=m\) is a
constant mass-conserving solution; \zcref{prop:closed} gives \(\sum\sigma_n=m(1-m)\) and \eqref{eq:closed} with
\(\dot\sigma=0\), that is \eqref{eq:recursion}; so \(\sigma=\rtil\). Relaxation: \(\rho_n(t)\to\rtil_n\) and both
masses equal \(m\), so \(\rho(t)\to\rtil\) in \(\ell^1_1\) by the discrete Scheffé argument.
\end{proof}

Jang and Tran obtained the stationary state for \(m\le1/2\) as the long-time limit of their solutions
(Lemma~5.7, using their Lemma~5.1 on viscosity solutions); for \(m\ge1\) none exists (their Lemma~5.6).

\section{Decay of the stationary state}\label{sec:asymptotics}

\emph{Attribution.} The singularity analysis of this section was found by GPT-6-Astra; it originally served a
continuation-in-\(m\) proof of positivity. Here positivity and the moment identities are supplied by
\zcref{thm:stationary}, so only the fixed-\(m\) analysis is needed; it was rechecked and condensed by the
final author.

Fix \(m\in(0,1)\), \(a=2m+1\), \(N=m(1-m)\), \(H(z)=\sum\rtil_nz^n\), \(P(z)=zH'(z)=\sum n\rtil_nz^n\), with common
radius \(R\). By \zcref{thm:stationary}, \(R\ge1\), \(P\) is strictly increasing on \([0,1]\) with
\(P(1)=m\), \(H(1)=N\). Multiplying \(K(\rtil)_n=0\) by \(2z^n\) and summing (absolutely convergent) gives, for
\(|z|<R\),
\begin{equation}\label{eq:gen}
(1-z)(aP-P^2)-2zN+(1+z)H=0,
\end{equation}
and, after differentiation and elimination of \(H\),
\begin{equation}\label{eq:P}
z(1-z^2)(a-2P)P'=2z(m-P)(P+1-m)-(1-z)^2P .
\end{equation}

\begin{lemma}\label{lem:one}
\(H\) and \(P\) extend analytically across \(z=1\); hence \(R>1\).
\end{lemma}
\begin{proof}
For \(0<z<1\) put \(r=(m-P)/(1-z)>0\). Substituting into \eqref{eq:P} gives
\[
r'=\frac{m-r+2z(z+2)r^2}{z(1+z)\bigl(1+2(1-z)r\bigr)} .
\]
Suppose \(r\) is unbounded near \(1\). For \(z\) close to \(1\) the numerator is positive once \(r\) exceeds a
fixed constant, so \(r\) eventually increases to \(+\infty\); since \((1-z)r=m-P\to0\), the denominator tends
to \(2\) and \(r'/r^2\to3\). Integrating \((1/r)'=-r'/r^2\) up to \(1\) gives \((1-z)r\to1/3\), a contradiction.
So \(r\) is bounded, \(r'\) is bounded, \(r\) has a limit \(r_1\), and the equation is analytic at
\((1,r_1)\) (denominator \(2\)); its analytic solution through that point coincides with \(r\). Thus
\(P=m-(1-z)r\) and, by \eqref{eq:gen}, \(H=N-\frac{1-z}{1+z}(N+aP-P^2)\) are analytic at \(1\). Pringsheim's
theorem \autocite[Theorem IV.6]{FlajoletSedgewick2009} (the coefficients are positive) excludes \(R=1\).
\end{proof}

\begin{lemma}\label{lem:turning}
\(R<\infty\), \(P(R^-)=p_*:=m+\frac12\), and near \(R\),
\(P(z)=p_*-b\sqrt{1-z/R}+O(1-z/R)\) with a convergent expansion in powers of \(\sqrt{1-z/R}\) and \(b\) as in
\zcref{thm:asymptotics}. The point \(R\) is the only singularity of \(H\) on \(|z|=R\).
\end{lemma}
\begin{proof}
On \((1,R)\), \(P>P(1)=m\). If \(P\) reached \(p_*\) at some \(z_0<R\), the left side of \eqref{eq:P} would
vanish there while the right side equals \(-\frac32z_0-(1-z_0)^2p_*<0\); so \(m<P<p_*\) on \((1,R)\). For
\(z\ge2\) in that interval, both sides of \eqref{eq:P} are negative and
\(P'\ge\frac{(z-1)P}{z(z+1)(a-2P)}\ge\frac P{3z}\) (as \(0<a-2P<1\)); if \(R=\infty\) this contradicts
\(P<p_*\). If \(P(R^-)<p_*\), \eqref{eq:P} is regular at \((R,P(R^-))\) and continues \(P\) analytically
across \(R\), contradicting Pringsheim; so \(P(R^-)=p_*\). Near this point the inverse function solves the
regular equation \(dz/dp=B(z)(a-2p)/D(z,p)\), \(B=z(1-z^2)\), \(D=2z(m-p)(p+1-m)-(1-z)^2p\), with
\(B(R)<0\), \(D(R,p_*)<0\); so \(z(p)\) is analytic at \(p_*\) with \(z'(p_*)=0\),
\(z''(p_*)=-2B(R)/D(R,p_*)<0\). Inverting this quadratic critical point gives the expansion, with
\(b^2=RD(R,p_*)/B(R)\), which is the stated formula. Finally \(\sum n\rtil_nR^n=p_*<\infty\), so \(H,P\) are
continuous on \(|z|\le R\). If \(|w|=R\), \(w\ne R\), strict positivity of \(\rtil_1,\rtil_2\) gives
\(|P(w)|<P(R)=a/2\), so \(a-2P(w)\neq0\) and \(w\ne1\); then \eqref{eq:gen} defines \(P\) analytically in
terms of \((z,H)\) near \((w,H(w),P(w))\), and the regular equation \(H'=P/z\) continues \(H\) analytically
across \(w\).
\end{proof}

\begin{proof}[of \zcref{thm:asymptotics}]
By \zcref{lem:turning} and compactness, \(H\) is analytic in a \(\Delta\)-domain at \(R\)
\autocite[Definition VI.1, Section VI.3]{FlajoletSedgewick2009}. With \(s=1-z/R\), \(dH/ds=-P/(1-s)\), so integrating the
expansion of \(P\) gives \(H=A(s)+s^{3/2}g(s)\) with \(A,g\) analytic near \(0\) and \(g(0)=\frac{2b}3\). Hence
\(H-A_0-A_1s-A_2s^2-\frac{2b}3s^{3/2}=O(s^{5/2})\) in the \(\Delta\)-domain, where \(A_0+A_1s+A_2s^2\) is a
polynomial in \(z\) and does not contribute to coefficients of index \(n>2\).
The transfer theorems \autocite[Theorems VI.1 and VI.3]{FlajoletSedgewick2009} give
\([z^n]H=\frac{2b}{3\,\Gamma(-3/2)}R^{-n}n^{-5/2}(1+O(n^{-1}))\), and \(\Gamma(-\frac32)=\frac43\sqrt\pi\).
\end{proof}

In particular \(\rtil\) has exponential moments \(\sum q^n\rtil_n<\infty\) for \(q<R(m)\). The slow decay for
\(m\) close to \(1\) (where \(R(m)\to1\)) is visible numerically: for \(m=0.75\), the mass carried by sizes
above \(3000\) is about \(2\times10^{-9}\).

\section{Remarks and limitations}

\begin{remark}[what is used where]
The restriction \(m\le1\) enters only through the face \(U=0\) of the invariant region
(\(\dot U=\frac12(P^2+m-m^2)\) there); for \(m>1\) Jang and Tran's Theorem~1.1 excludes global mass-conserving
solutions, and \zcref{prop:closed} shows directly that \(n_0\) would become negative. The restriction \(m<1\)
enters only in the uniform-in-time estimate \zcref{lem:slope}(b), which is needed for the stationary state.
\end{remark}

\begin{remark}[moments]
\zcref{thm:evolution} does not assert that \(m_2(\rho(t))\) is finite for \(t>0\) when \(m_2(\rho^0)=\infty\);
it is not needed. The solution is classical in Jang and Tran's sense (each component is differentiable and
satisfies \eqref{eq:system}).
\end{remark}

\begin{remark}[relation to the literature]
Tran and Van \autocite{TranVan2022}, Jang and Tran, and Van construct solutions as vanishing-viscosity limits of viscosity
solutions and then prove a Bernstein (positivity) property, which requires regularity of the
Hamilton--Jacobi solution at the boundary \(x=0\). Here positivity comes from the finite truncations, the
characteristic solution is used only as a lower barrier, and its boundary behaviour reduces to the ODE
inequality of \zcref{lem:slope}. The invariant \(W\) is the discrete form of Van's half-slope invariant
\(2M-xM_x\ge0\), \(M=mx-F\) \autocite{Van2026}.
\end{remark}

\begin{remark}[an alternative stationary argument]
GPT-6-Astra also sketched a purely stationary proof of positivity of \eqref{eq:recursion}, by continuation in
\(m\) from Jang and Tran's range \(m\le1/2\), based on a parameter-uniform version of
\zcref{lem:one,lem:turning} and uniform singularity transfer. Its uniformity step (a common
\(\Delta\)-domain and uniform transfer near the moving turning point \(R(m)\)) has not been checked to the
standard of this paper, and it is not used; the draft is kept outside the deliverables.
\end{remark}

\section{Computer checks and independent verification}

The proofs do not depend on numerical computation. Scripts in \texttt{scripts/} (each prints PASS/FAIL):
\texttt{check\_evolution\_identities.py} verifies exactly (sympy) every identity of
Sections~\ref{sec:HJ}--\ref{sec:existence}: \(1/\tau=\frac12+\frac1{e^x-1}\), the transform equations for the
full and truncated systems (\eqref{eq:super} with \(E_L\)), the characteristic equations, \eqref{eq:W},
\eqref{eq:WD}, \eqref{eq:lambda}, the strip condition, the Riccati equation, \eqref{eq:Gamma}, the identity of
\zcref{lem:initial} for \(j\le60\) and the elementary inequalities.
\texttt{check\_identities.py} (GPT-6-Astra) verifies the stationary identities \eqref{eq:gen}, \eqref{eq:P},
the \(r\)-equation, the turning-point quantities and \(b^2\); \texttt{check\_coefficients\_arb.py} encloses the
first 256 coefficients of \eqref{eq:recursion} at six masses in Arb ball arithmetic (finite consistency check).
\texttt{numerics\_evolution.py} is floating-point evidence only: invariants and barrier along characteristics
for monomer and infinite-second-moment data, the comparison \(F^L\ge F\), agreement of the triangular and
truncated systems, and the long-time limit. Its complete run is reported in \texttt{README.md}.

\paragraph{Independent verification.} The points singled out for checking (the no-crossing and
global-coverage argument, \zcref{lem:barrier} and \zcref{prop:global}; the uniformity in \zcref{lem:slope}(a)
over characteristics arriving near \(x=0\); the identity \eqref{eq:super} and the sign of \(E_L\); the
tightness step in the coagulation loss; the use of \zcref{lem:slope}(b) to prevent mass escape as
\(t\to\infty\); the singularity analysis) were re-derived by an independent verifier, who found no
mathematical gap. The verifier re-ran all scripts and added independent checks (different integrators,
infinite-second-moment data at \(m=0.9\) and \(m=1\), Arb enclosures of \eqref{eq:recursion}, an independent
computation of \(R(m)\) and of the constant of \zcref{thm:asymptotics}) in \texttt{verification/}; the
report, with the list of corrections, is \texttt{verification.md}.

\begingroup
\emergencystretch=2em
\paperreferences
\endgroup
\end{document}
